\documentclass[11pt,a4paper]{article}

\usepackage[utf8]{inputenc}
\usepackage[T1]{fontenc}
\usepackage[french]{babel}
\usepackage[a4paper,margin=2.2cm]{geometry}
\usepackage{amsmath,amssymb,amsthm,mathtools}
\usepackage{lmodern}
\usepackage{enumitem}
\usepackage{hyperref}
\usepackage{xcolor}
\usepackage{fancyhdr}
\usepackage{tikz}

\hypersetup{
    colorlinks=true,
    linkcolor=black,
    urlcolor=blue
}

\setlength{\parindent}{0pt}
\setlength{\parskip}{0.6em}

\pagestyle{fancy}
\fancyhf{}
\fancyhead[L]{Fiche d’exercices --- Algèbre linéaire Chapitre 4}
\fancyhead[R]{Familles libres, génératrices et bases}
\fancyfoot[C]{\thepage}

\begin{document}

\begin{center}
    {\LARGE \textbf{Fiche d’exercices --- Chapitre 4}}\\[0.4em]
    {\Large \textbf{Familles libres, génératrices et bases}}\\[0.6em]
    Algèbre linéaire
\end{center}

\textbf{AIDE de lecture type d’exo :} APP = application directe, METH = méthode, DEM = démonstration, PREP = réflexion plus poussée. Les étoiles représentent un gradient de difficulté de 1 à 5.

\tableofcontents

\newpage

\section*{Exercices d’application directe}
\addcontentsline{toc}{section}{Exercices d’application directe}

\textbf{Exercice 1 [APP $\star$]}\\
\textbf{Libre, génératrice, base dans \(\mathbb{R}^2\)}\\
Pour chacune des familles suivantes de \(\mathbb{R}^2\), dire si elle est libre, génératrice, et si c’est une base :
\begin{enumerate}[label=\alph*)]
    \item \(\big((1,0),(0,1)\big)\),
    \item \(\big((1,0),(2,0)\big)\),
    \item \(\big((1,1),(1,-1)\big)\),
    \item \(\big((1,2)\big)\),
    \item \(\big((1,0),(0,1),(1,1)\big)\).
\end{enumerate}

\textbf{Exercice 2 [APP $\star$]}\\
\textbf{Famille à un vecteur}\\
Dire, selon les cas, si la famille \((u)\) est libre :
\begin{enumerate}[label=\alph*)]
    \item dans \(\mathbb{R}^2\), avec \(u=(1,2)\),
    \item dans \(\mathbb{R}^2\), avec \(u=(0,0)\),
    \item dans \(\mathbb{R}[X]\), avec \(u=X^2+1\),
    \item dans \(\mathcal F(\mathbb{R},\mathbb{R})\), avec \(u=0\).
\end{enumerate}

\textbf{Exercice 3 [APP $\star\star$]}\\
\textbf{Vecteur redondant}\\
Dans chacune des familles suivantes, déterminer si l’un des vecteurs est combinaison linéaire des autres :
\begin{enumerate}[label=\alph*)]
    \item \(\big((1,0),(0,1),(1,1)\big)\),
    \item \(\big((1,2),(2,4)\big)\),
    \item \(\big((1,0,0),(0,1,0),(1,1,0)\big)\),
    \item \(\big(1,X,X+1\big)\) dans \(\mathbb{R}[X]\).
\end{enumerate}

\textbf{Exercice 4 [APP $\star\star$]}\\
\textbf{Première étude de liberté}\\
Étudier la liberté des familles suivantes :
\begin{enumerate}[label=\alph*)]
    \item \(\big((1,2),(2,3)\big)\) dans \(\mathbb{R}^2\),
    \item \(\big((1,2),(2,4)\big)\) dans \(\mathbb{R}^2\),
    \item \(\big((1,0,1),(0,1,1),(1,1,0)\big)\) dans \(\mathbb{R}^3\),
    \item \(\big(1,X,X^2\big)\) dans \(\mathbb{R}[X]\).
\end{enumerate}

\textbf{Exercice 5 [APP $\star\star$]}\\
\textbf{Première étude du caractère générateur}\\
Déterminer si les familles suivantes sont génératrices de l’espace indiqué :
\begin{enumerate}[label=\alph*)]
    \item \(\big((1,0),(0,1)\big)\) dans \(\mathbb{R}^2\),
    \item \(\big((1,1),(1,-1)\big)\) dans \(\mathbb{R}^2\),
    \item \(\big((1,0,0),(0,1,0)\big)\) dans \(\mathbb{R}^3\),
    \item \(\big(1,X,X^2\big)\) dans \(\mathbb{R}_2[X]\) (polynômes de degré \(\le 2\)).
\end{enumerate}

\textbf{Exercice 6 [APP $\star\star$]}\\
\textbf{Coordonnées dans la base canonique}\\
Donner les coordonnées des vecteurs suivants dans la base canonique :
\begin{enumerate}[label=\alph*)]
    \item \((3,-2)\) dans \(\mathbb{R}^2\),
    \item \((1,0,4)\) dans \(\mathbb{R}^3\),
    \item \(2-3X+X^2\) dans la base \((1,X,X^2)\) de \(\mathbb{R}_2[X]\).
\end{enumerate}

\textbf{Exercice 7 [APP $\star\star$]}\\
\textbf{Coordonnées dans une autre base}\\
Dans la base \(\mathcal B=\big((1,1),(1,-1)\big)\) de \(\mathbb{R}^2\), déterminer les coordonnées de :
\begin{enumerate}[label=\alph*)]
    \item \((3,1)\),
    \item \((5,-1)\),
    \item \((0,2)\),
    \item \((a,b)\) avec \(a,b\in\mathbb{R}\).
\end{enumerate}

\textbf{Exercice 8 [APP $\star\star$]}\\
\textbf{Sous-familles et sur-familles}\\
Pour chacune des affirmations suivantes, dire si elle est vraie ou fausse :
\begin{enumerate}[label=\alph*)]
    \item toute sous-famille d’une famille libre est libre ;
    \item toute sous-famille d’une famille génératrice est génératrice ;
    \item toute sur-famille d’une famille génératrice est génératrice ;
    \item toute sur-famille d’une famille libre est libre.
\end{enumerate}

\textbf{Exercice 9 [APP $\star\star$]}\\
\textbf{Base canonique}\\
Montrer que la famille canonique est une base de :
\begin{enumerate}[label=\alph*)]
    \item \(\mathbb{R}^2\),
    \item \(\mathbb{R}^3\),
    \item \(\mathbb{R}^n\).
\end{enumerate}

\textbf{Exercice 10 [APP $\star\star$]}\\
\textbf{Polynômes}\\
Dans \(\mathbb{R}_2[X]\), étudier les familles suivantes :
\begin{enumerate}[label=\alph*)]
    \item \((1,X,X^2)\),
    \item \((1,X,1+X)\),
    \item \((X,X^2)\),
    \item \((1+X, X+X^2, 1+X^2)\).
\end{enumerate}

\newpage

\section*{Exercices de méthode}
\addcontentsline{toc}{section}{Exercices de méthode}

\textbf{Exercice 11 [METH $\star\star$]}\\
\textbf{Méthode générale pour la liberté}\\
Étudier la liberté de la famille
\[
\big((1,0,1),(0,1,1),(1,1,0)\big)
\]
dans \(\mathbb{R}^3\) en suivant précisément la méthode du cours :
\[
\lambda_1u_1+\lambda_2u_2+\lambda_3u_3=0.
\]

\textbf{Exercice 12 [METH $\star\star$]}\\
\textbf{Méthode générale pour le caractère générateur}\\
Montrer que la famille
\[
\big((1,0,1),(0,1,1),(1,1,0)\big)
\]
engendre \(\mathbb{R}^3\) en partant d’un vecteur arbitraire \((x,y,z)\).

\textbf{Exercice 13 [METH $\star\star\star$]}\\
\textbf{Base ou non ?}\\
Déterminer si les familles suivantes sont des bases de l’espace indiqué :
\begin{enumerate}[label=\alph*)]
    \item \(\big((1,2),(2,3)\big)\) de \(\mathbb{R}^2\),
    \item \(\big((1,0,1),(0,1,1),(1,1,0)\big)\) de \(\mathbb{R}^3\),
    \item \(\big(1,X,X^2\big)\) de \(\mathbb{R}_2[X]\),
    \item \(\big(1,\mathrm{id}\big)\) de l’espace des polynômes de degré \(\le 1\).
\end{enumerate}

\textbf{Exercice 14 [METH $\star\star\star$]}\\
\textbf{Base d’un sous-espace de \(\mathbb{R}^3\)}\\
Déterminer une base de
\[
F=\{(x,y,z)\in\mathbb{R}^3\mid x+y-z=0\}.
\]
On demandera une rédaction par analyse-synthèse.

\textbf{Exercice 15 [METH $\star\star\star$]}\\
\textbf{Autre base d’un sous-espace}\\
Déterminer une base de
\[
F=\{(x,y,z)\in\mathbb{R}^3\mid x-y+z=0\}.
\]

\textbf{Exercice 16 [METH $\star\star\star$]}\\
\textbf{Extraction d’une base à partir d’une famille génératrice}\\
Dans \(\mathbb{R}^3\), on considère la famille
\[
\mathcal F=\big((1,0,0),(0,1,0),(1,1,0),(0,0,1)\big).
\]
\begin{enumerate}[label=\alph*)]
    \item Montrer qu’elle est génératrice de \(\mathbb{R}^3\).
    \item Montrer qu’elle est liée.
    \item En extraire une base de \(\mathbb{R}^3\).
\end{enumerate}

\textbf{Exercice 17 [METH $\star\star\star$]}\\
\textbf{Deux écritures d’un même vecteur}\\
Soit \(\mathcal B=\big((1,1),(1,-1)\big)\), base de \(\mathbb{R}^2\).\\
On considère un vecteur \(u\in\mathbb{R}^2\) admettant deux écritures :
\[
u=\lambda(1,1)+\mu(1,-1)=\lambda'(1,1)+\mu'(1,-1).
\]
Montrer que \(\lambda=\lambda'\) et \(\mu=\mu'\).

\textbf{Exercice 18 [METH $\star\star\star$]}\\
\textbf{Déterminer les coordonnées dans une base non canonique}\\
Dans \(\mathbb{R}^3\), on considère la famille
\[
\mathcal B=\big((1,0,0),(1,1,0),(1,1,1)\big).
\]
\begin{enumerate}[label=\alph*)]
    \item Montrer que \(\mathcal B\) est une base de \(\mathbb{R}^3\).
    \item Déterminer les coordonnées de \((x,y,z)\) dans cette base.
    \item En déduire les coordonnées de \((2,-1,3)\).
\end{enumerate}

\textbf{Exercice 19 [METH $\star\star\star$]}\\
\textbf{Famille génératrice avec paramètre}\\
Dans \(\mathbb{R}^2\), on considère la famille
\[
\mathcal F_a=\big((1,0),(a,1)\big),\qquad a\in\mathbb{R}.
\]
Pour quelles valeurs de \(a\) la famille \(\mathcal F_a\) est-elle :
\begin{enumerate}[label=\alph*)]
    \item libre ?
    \item génératrice ?
    \item une base de \(\mathbb{R}^2\) ?
\end{enumerate}

\textbf{Exercice 20 [METH $\star\star\star$]}\\
\textbf{Famille à trois vecteurs dans \(\mathbb{R}^2\)}\\
Dans \(\mathbb{R}^2\), on considère
\[
u_1=(1,0),\qquad u_2=(0,1),\qquad u_3=(a,b).
\]
\begin{enumerate}[label=\alph*)]
    \item La famille \((u_1,u_2,u_3)\) est-elle libre ?
    \item Est-elle génératrice ?
    \item Dans quels cas contient-elle un vecteur redondant évident ?
\end{enumerate}

\textbf{Exercice 21 [METH $\star\star\star\star$]}\\
\textbf{Polynômes et base adaptée}\\
Dans \(\mathbb{R}_2[X]\), on considère la famille
\[
\mathcal F=\big(1,\;1+X,\;1+X+X^2\big).
\]
\begin{enumerate}[label=\alph*)]
    \item Montrer que \(\mathcal F\) est une base de \(\mathbb{R}_2[X]\).
    \item Déterminer les coordonnées de \(P(X)=2-X+3X^2\) dans cette base.
\end{enumerate}

\textbf{Exercice 22 [METH $\star\star\star\star$]}\\
\textbf{Fonctions affines}\\
Dans l’espace \(E\) des fonctions polynomiales de degré \(\le 1\), on considère la famille
\[
\mathcal F=(1,\mathrm{id}).
\]
\begin{enumerate}[label=\alph*)]
    \item Montrer que \(\mathcal F\) est libre.
    \item Montrer qu’elle engendre \(E\).
    \item En déduire que \(\mathcal F\) est une base de \(E\).
    \item Donner les coordonnées de la fonction \(x\mapsto 3x-2\) dans cette base.
\end{enumerate}

\newpage

\section*{Exercices de démonstration}
\addcontentsline{toc}{section}{Exercices de démonstration}

\textbf{Exercice 23 [DEM $\star\star$]}\\
\textbf{Caractérisation d’une famille liée}\\
Montrer qu’une famille finie \((u_1,\dots,u_p)\) est liée si et seulement s’il existe des scalaires non tous nuls \(\lambda_1,\dots,\lambda_p\) tels que
\[
\lambda_1u_1+\cdots+\lambda_pu_p=0.
\]

\textbf{Exercice 24 [DEM $\star\star$]}\\
\textbf{Famille contenant le vecteur nul}\\
Montrer que toute famille contenant le vecteur nul est liée.

\textbf{Exercice 25 [DEM $\star\star$]}\\
\textbf{Deux vecteurs égaux}\\
Montrer que si une famille contient deux vecteurs égaux, alors elle est liée.

\textbf{Exercice 26 [DEM $\star\star\star$]}\\
\textbf{Vecteur combinaison linéaire des autres}\\
Montrer que si un vecteur d’une famille est combinaison linéaire des autres, alors la famille est liée.

\textbf{Exercice 27 [DEM $\star\star\star$]}\\
\textbf{Caractérisation opérationnelle d’une famille liée}\\
Montrer qu’une famille finie est liée si et seulement si l’un de ses vecteurs est combinaison linéaire des autres.

\textbf{Exercice 28 [DEM $\star\star\star$]}\\
\textbf{Sous-famille d’une famille libre}\\
Montrer que toute sous-famille d’une famille libre est libre.

\textbf{Exercice 29 [DEM $\star\star\star$]}\\
\textbf{Sur-famille d’une famille génératrice}\\
Montrer que toute famille contenant une famille génératrice est génératrice.

\textbf{Exercice 30 [DEM $\star\star\star$]}\\
\textbf{Liberté et unicité d’écriture}\\
Montrer qu’une famille finie \((u_1,\dots,u_p)\) est libre si et seulement si tout vecteur de \(\mathrm{Vect}(u_1,\dots,u_p)\) admet une écriture unique comme combinaison linéaire de cette famille.

\textbf{Exercice 31 [DEM $\star\star\star$]}\\
\textbf{Une famille libre est une base de son espace engendré}\\
Montrer que toute famille libre est une base du sous-espace qu’elle engendre.

\textbf{Exercice 32 [DEM $\star\star\star$]}\\
\textbf{Base canonique}\\
Montrer que la base canonique de \(K^n\) est bien une base de \(K^n\).

\textbf{Exercice 33 [DEM $\star\star\star\star$]}\\
\textbf{Extraction d’une base}\\
Montrer que de toute famille génératrice finie d’un espace vectoriel, on peut extraire une base.

\textbf{Exercice 34 [DEM $\star\star\star\star$]}\\
\textbf{Polynômes de degré \(\le 2\)}\\
Montrer que la famille \((1,X,X^2)\) est une base de \(\mathbb{R}_2[X]\).

\textbf{Exercice 35 [DEM $\star\star\star\star$]}\\
\textbf{Fonctions \(1\) et \(\mathrm{id}\)}\\
Montrer que la famille \((1,\mathrm{id})\) est libre dans l’espace des fonctions de \(\mathbb{R}\) dans \(\mathbb{R}\).

\newpage

\section*{Exercices de réflexion}
\addcontentsline{toc}{section}{Exercices de réflexion}

\textbf{Exercice 36 [PREP $\star\star\star$]}\\
\textbf{Libre mais non génératrice ?}\\
Donner un exemple :
\begin{enumerate}[label=\alph*)]
    \item d’une famille libre mais non génératrice de \(\mathbb{R}^3\),
    \item d’une famille génératrice mais liée de \(\mathbb{R}^2\),
    \item d’une famille qui ne soit ni libre ni génératrice dans \(\mathbb{R}^3\),
    \item d’une base de \(\mathbb{R}^2\) différente de la base canonique.
\end{enumerate}

\textbf{Exercice 37 [PREP $\star\star\star$]}\\
\textbf{Trois vecteurs dans \(\mathbb{R}^2\)}\\
Peut-on trouver trois vecteurs de \(\mathbb{R}^2\) qui soient libres ?\\
Peut-on trouver trois vecteurs de \(\mathbb{R}^2\) qui soient générateurs ?\\
Donner des justifications précises.

\textbf{Exercice 38 [PREP $\star\star\star$]}\\
\textbf{Deux vecteurs dans \(\mathbb{R}^3\)}\\
Peut-on trouver deux vecteurs de \(\mathbb{R}^3\) qui soient générateurs de \(\mathbb{R}^3\) ?\\
Peut-on trouver deux vecteurs de \(\mathbb{R}^3\) qui soient libres ?\\
Justifier.

\textbf{Exercice 39 [PREP $\star\star\star\star$]}\\
\textbf{Famille libre et relation cachée}\\
Dans \(\mathbb{R}^3\), on considère
\[
u_1=(1,1,0),\qquad u_2=(1,0,1),\qquad u_3=(0,1,1).
\]
\begin{enumerate}[label=\alph*)]
    \item Étudier la liberté de \((u_1,u_2,u_3)\).
    \item Étudier si cette famille est génératrice de \(\mathbb{R}^3\).
    \item Déterminer si c’est une base.
\end{enumerate}

\textbf{Exercice 40 [PREP $\star\star\star\star$]}\\
\textbf{Description explicite d’un sous-espace engendré}\\
Décrire explicitement
\[
\mathrm{Vect}\big((1,0,-1),(0,1,-1)\big)\subset\mathbb{R}^3.
\]
Puis montrer que c’est exactement le sous-espace
\[
F=\{(x,y,z)\in\mathbb{R}^3\mid x+y+z=0\}.
\]

\textbf{Exercice 41 [PREP $\star\star\star\star$]}\\
\textbf{Coordonnées et paramètre}\\
Dans \(\mathbb{R}^2\), on considère la famille
\[
\mathcal B_a=\big((1,1),(1,a)\big).
\]
\begin{enumerate}[label=\alph*)]
    \item Pour quelles valeurs de \(a\) cette famille est-elle une base de \(\mathbb{R}^2\) ?
    \item Lorsque c’est une base, déterminer les coordonnées de \((x,y)\) dans \(\mathcal B_a\).
\end{enumerate}

\textbf{Exercice 42 [PREP $\star\star\star\star$]}\\
\textbf{Coordonnées dans une base de polynômes}\\
Dans \(\mathbb{R}_2[X]\), on considère la base
\[
\mathcal B=\big(1,\;X-1,\;(X-1)^2\big).
\]
\begin{enumerate}[label=\alph*)]
    \item Montrer que \(\mathcal B\) est une base de \(\mathbb{R}_2[X]\).
    \item Déterminer les coordonnées de \(P(X)=X^2+X+1\) dans cette base.
\end{enumerate}

\textbf{Exercice 43 [PREP $\star\star\star\star$]}\\
\textbf{Une famille de polynômes dépendant d’un paramètre}\\
Dans \(\mathbb{R}_2[X]\), on considère
\[
\mathcal F_a=(1,\;X+a,\;X^2).
\]
Étudier, selon \(a\in\mathbb{R}\), si cette famille est libre, génératrice, puis base de \(\mathbb{R}_2[X]\).

\textbf{Exercice 44 [PREP $\star\star\star\star$]}\\
\textbf{Une famille de fonctions}\\
Dans l’espace des fonctions de \(\mathbb{R}\) dans \(\mathbb{R}\), on considère
\[
f_1(x)=1,\qquad f_2(x)=x,\qquad f_3(x)=x+1.
\]
\begin{enumerate}[label=\alph*)]
    \item Étudier la liberté de \((f_1,f_2,f_3)\).
    \item Déterminer \(\mathrm{Vect}(f_1,f_2,f_3)\).
\end{enumerate}

\textbf{Exercice 45 [PREP $\star\star\star\star$]}\\
\textbf{Une base d’un plan de \(\mathbb{R}^3\)}\\
On considère
\[
F=\{(x,y,z)\in\mathbb{R}^3\mid 2x-y+z=0\}.
\]
\begin{enumerate}[label=\alph*)]
    \item Déterminer une base de \(F\).
    \item Le vecteur \((1,1,-1)\) appartient-il à \(F\) ?
    \item Donner les coordonnées de tout vecteur de \(F\) dans la base trouvée.
\end{enumerate}

\textbf{Exercice 46 [PREP $\star\star\star\star\star$]}\\
\textbf{Famille libre, génératrice, base : panorama complet}\\
Dans \(\mathbb{R}^3\), on considère les familles suivantes :
\[
\mathcal F_1=\big((1,0,0),(0,1,0)\big),\qquad
\mathcal F_2=\big((1,0,0),(0,1,0),(1,1,0)\big),
\]
\[
\mathcal F_3=\big((1,0,0),(0,1,0),(0,0,1)\big),\qquad
\mathcal F_4=\big((1,1,0),(1,0,1),(0,1,1)\big).
\]
Pour chacune :
\begin{enumerate}[label=\alph*)]
    \item déterminer le sous-espace engendré ;
    \item dire si la famille est libre ;
    \item dire si elle est génératrice de \(\mathbb{R}^3\) ;
    \item dire si c’est une base de \(\mathbb{R}^3\).
\end{enumerate}

\textbf{Exercice 47 [PREP $\star\star\star\star\star$]}\\
\textbf{Extraire intelligemment une base}\\
Dans \(\mathbb{R}^4\), on considère
\[
u_1=(1,0,0,0),\quad
u_2=(0,1,0,0),\quad
u_3=(1,1,0,0),\quad
u_4=(0,0,1,0),\quad
u_5=(0,0,0,1).
\]
\begin{enumerate}[label=\alph*)]
    \item Montrer que la famille \((u_1,u_2,u_3,u_4,u_5)\) est génératrice de \(\mathbb{R}^4\).
    \item Montrer qu’elle est liée.
    \item En extraire une base de \(\mathbb{R}^4\).
    \item Toutes les extractions possibles donnent-elles la même base ?
\end{enumerate}

\textbf{Exercice 48 [PREP $\star\star\star\star\star$]}\\
\textbf{Un problème de coordonnées complet}\\
Dans \(\mathbb{R}^3\), on considère
\[
\mathcal B=\big((1,0,1),(0,1,1),(1,1,0)\big).
\]
\begin{enumerate}[label=\alph*)]
    \item Montrer que \(\mathcal B\) est une base de \(\mathbb{R}^3\).
    \item Déterminer les coordonnées de \((x,y,z)\) dans \(\mathcal B\).
    \item En déduire les coordonnées de \((1,2,3)\).
    \item Déterminer le vecteur dont les coordonnées dans \(\mathcal B\) sont \((2,-1,4)\).
\end{enumerate}

\textbf{Exercice 49 [PREP $\star\star\star\star\star$]}\\
\textbf{Analyse-synthèse dans \(\mathbb{R}^3\)}\\
Déterminer une base du sous-espace
\[
F=\{(x,y,z)\in\mathbb{R}^3\mid x+y+z=0\}
\]
puis, pour le vecteur \(u=(1,-2,1)\), déterminer ses coordonnées dans cette base.

\textbf{Exercice 50 [PREP $\star\star\star\star\star$]}\\
\textbf{Problème de synthèse final}\\
Pour chacune des affirmations suivantes :
\begin{enumerate}[label=\alph*)]
    \item dire si elle est vraie ou fausse ;
    \item si elle est vraie, la démontrer ;
    \item si elle est fausse, donner un contre-exemple.
\end{enumerate}

\begin{enumerate}[label=\arabic*)]
    \item Toute famille libre est génératrice.
    \item Toute famille génératrice est libre.
    \item Toute base est une famille libre.
    \item Toute base est une famille génératrice.
    \item Toute famille libre est une base de son espace engendré.
    \item Toute famille contenant une base est une base.
    \item Toute sous-famille d’une base est libre.
    \item Toute sur-famille d’une base est génératrice.
\end{enumerate}

\newpage

\section*{Pour aller plus loin}
\addcontentsline{toc}{section}{Pour aller plus loin}

\textbf{Exercice 51 [PREP $\star\star\star\star\star$]}\\
\textbf{Famille libre de fonctions}\\
Dans l’espace des fonctions de \(\mathbb{R}\) dans \(\mathbb{R}\), montrer que la famille
\[
(1,\mathrm{id},\mathrm{id}^2)
\]
est libre, où \(\mathrm{id}(x)=x\) et \(\mathrm{id}^2(x)=x^2\).

\textbf{Exercice 52 [PREP $\star\star\star\star\star$]}\\
\textbf{Polynômes et coordonnées abstraites}\\
Dans \(\mathbb{R}_2[X]\), on considère la base
\[
\mathcal B=\big(1,\;1+X,\;1+X+X^2\big).
\]
\begin{enumerate}[label=\alph*)]
    \item Déterminer les coordonnées de \(1\), \(X\) et \(X^2\) dans \(\mathcal B\).
    \item Retrouver ensuite les coordonnées d’un polynôme général \(a+bX+cX^2\) dans \(\mathcal B\).
\end{enumerate}

\textbf{Exercice 53 [PREP $\star\star\star\star\star$]}\\
\textbf{Base d’un sous-espace de polynômes}\\
Soit
\[
F=\{P\in\mathbb{R}_3[X]\mid P(1)=0\}.
\]
\begin{enumerate}[label=\alph*)]
    \item Montrer que \(F\) est un sous-espace vectoriel de \(\mathbb{R}_3[X]\).
    \item Montrer que tout polynôme de \(F\) est multiple de \(X-1\).
    \item En déduire une famille génératrice simple de \(F\).
    \item Montrer que cette famille est une base de \(F\).
\end{enumerate}

\textbf{Exercice 54 [PREP $\star\star\star\star\star$]}\\
\textbf{Familles avec paramètre dans \(\mathbb{R}^3\)}\\
Dans \(\mathbb{R}^3\), on considère
\[
u_1=(1,0,0),\qquad u_2=(0,1,0),\qquad u_3=(a,b,1).
\]
\begin{enumerate}[label=\alph*)]
    \item Étudier, selon \(a,b\in\mathbb{R}\), si \((u_1,u_2,u_3)\) est libre.
    \item Même question pour le caractère générateur.
    \item En déduire pour quelles valeurs de \(a,b\) cette famille est une base de \(\mathbb{R}^3\).
\end{enumerate}

\textbf{Exercice 55 [PREP $\star\star\star\star\star$]}\\
\textbf{Synthèse finale longue}\\
Dans \(\mathbb{R}^3\), on considère
\[
u_1=(1,0,1),\qquad u_2=(1,1,0),\qquad u_3=(0,1,1),\qquad u_4=(2,1,1).
\]
\begin{enumerate}[label=\alph*)]
    \item Étudier la liberté de \((u_1,u_2,u_3)\).
    \item Montrer que \((u_1,u_2,u_3)\) est une base de \(\mathbb{R}^3\).
    \item Déterminer les coordonnées de \(u_4\) dans cette base.
    \item En déduire que la famille \((u_1,u_2,u_3,u_4)\) est liée.
    \item Extraire de cette famille une base de \(\mathbb{R}^3\).
\end{enumerate}

\end{document}